% ---------------------------------------------------------------------
\section{Breit 系与领头区域；两条途径}\label{sec:disall-regions}

先固定参考系与 $D$ 维投影算符（$W^{\mu\nu}$、$T^{\mu\nu}$ 见\cref{def:dis-kinematics,def:dis-compton}）。

\begin{definition}[Breit 系与投影算符]\label{def:disall-breit}
设靶强子动量 $P$（质量 $M$）沿 $+z$ 方向，光锥分量记作 $(v^+,v^-,\mathbf v_T)$。取
\begin{equation}\label{eq:disall-breit}
P=\Big(P^+,\frac{M^2}{2P^+},\mathbf 0_T\Big),\qquad q=\Big(-x_NP^+,\frac{Q^2}{2x_NP^+},\mathbf 0_T\Big),
\end{equation}
其中 $x_N\equiv2x/(1+\rho)$，$\rho\equiv\sqrt{1+4x^2M^2/Q^2}$。于是 $q^2=-Q^2$ 与 $2P\cdot q=Q^2/x$ 精确成立，$x_N=x-x^3M^2/Q^2+\order{M^4/Q^4}$（Nachtmann 变量）。$P^+$ 可任意选取（沿 $z$ 轴的推动）；$P^+=Q/(\sqrt2\,x_N)$ 时 $q=(0,\mathbf 0_T,-Q)$（笛卡儿分量），即 Breit（“砖墙”）系本义。
在 $D=4-2\varepsilon$ 维中以同样的张量分解定义 $W^{\mu\nu}$，则对无质量靶
\begin{equation}\label{eq:disall-proj}
\frac{F_L}{x}=\frac{8x^2}{Q^2}\,P_\mu P_\nu W^{\mu\nu},\qquad
\frac{F_2}{x}=\frac1{1-\varepsilon}\Big[-g_{\mu\nu}+(3-2\varepsilon)\frac{4x^2}{Q^2}P_\mu P_\nu\Big]W^{\mu\nu};   % VF-DIS-2
\end{equation}
对有质量靶，以 $F_L\equiv\rho^2F_2-2xF_1$ 定义纵向结构函数，
\begin{equation}\label{eq:disall-proj-massive}
\frac{F_L}{x}=\frac{8x^2}{Q^2\rho^2}\,P_\mu P_\nu W^{\mu\nu},\qquad
\frac{\rho^2F_2}{x}=\frac{2(-g_{\mu\nu}W^{\mu\nu})+(D-1)F_L/x}{D-2}.   % VF-DIS-2
\end{equation}
对动量为 $p$ 的部分子靶，以 $p$ 代替 $P$、以 $z\equiv Q^2/(2p\cdot q)$ 代替 $x$，得到部分子结构函数 $\hat F_{2,L}$。
\uses{def:dis-kinematics}
\end{definition}

\noindent
投影算符来自直接缩并：记 $\tilde P=P-(P\cdot q/q^2)q$，$\tilde P^2=Q^2\rho^2/(4x^2)$，则 $-g_{\mu\nu}W^{\mu\nu}=(D-1)F_1-\tilde P^2F_2/(P\cdot q)$，
$P_\mu P_\nu W^{\mu\nu}=\frac{Q^2\rho^2}{8x^3}(\rho^2F_2-2xF_1)$。常用的 $q^+=-xP^+$ 给出 $2P\cdot q=Q^2/x-xM^2$，只在领头幂次正确。对无质量靶，被击夸克的虚度为
\begin{equation}\label{eq:disall-current-virtuality}
(\xi P+q)^2=2\xi P\cdot q-Q^2=Q^2(\xi-x)/x;
\end{equation}
领头阶 $\xi=x$，被击夸克沿 PDF Wilson 线的方向 $n_B=(0^+,1^-,\mathbf 0_T)$ 飞出。对固定 $x<1$ 的单个割图，领头钳夹面由靶共线子图、硬子图、沿 $n_B$ 的“流喷注”与软子图组成（\cref{fig:reduced-dis}）；
在非物理区的\emph{前向振幅}中后两者都不出现。

\paragraph{非物理区中的领头区域。}
对 $|\omega|<1$ 取参考系 $q^+=0$、$|\mathbf q_T|=Q$，于是 $2P\cdot q=2P^+q^-=\omega Q^2$，这一参数化对复 $\omega$ 也成立。
设 $k=\xi P+\kappa$ 是从靶共线子图流入硬部分的动量，$\kappa^+=0$，$|\boldsymbol\kappa_T|\lesssim\lambda Q$，$|\kappa^-|\lesssim\lambda^2Q^2/P^+$；由 $\kappa\cdot q=-\boldsymbol\kappa_T\cdot\mathbf q_T$，任何携带 $q$ 的线满足
\begin{equation}\label{eq:disall-unphysical-offshell}
(k+q)^2=Q^2(\xi\omega-1)+2\kappa\cdot q+k^2=Q^2(\xi\omega-1)+\order{\lambda Q^2},\qquad |\xi\omega-1|\ge1-|\omega|\quad(|\xi|\le1).
\end{equation}
这里 $|\xi|\le1$，因为份额 $|\xi|>1$ 时靶共线子图的 $k^-$ 极点全在实轴同一侧，不形成钳夹（与\cref{thm:pdf-support-pt}同一论证）。所以对 $|\omega|\le\rho<1$，携带 $q$ 的线离壳至少 $Q^2(1-\rho)$，在每个区域中都是硬线，附着在它们上的软线随之进入硬子图（附着在硬子图上的软线被幂次压低，\cref{ch:ir}），
只附着在靶共线子图上的软线属于共线子图内部。领头区域因此只由靶共线子图 $J$ 与硬子图 $H$ 组成，二者由两条物理极化的部分子外加任意多条标量（纵向）极化的 K 胶子相连；
K 胶子对 $H$ 的附着之和经 Ward 恒等式成为 PDF 的 Wilson 线（一般情形是\cref{lem:method-ward-multi}的内容，全阶论证，[A]），其种子见\cref{fig:disall-seed}。

%%% BEGIN fig:disall-seed
\begin{figure}[htbp]
\centering
\begin{tikzpicture}[fd,
  fb lab/.style={font=\footnotesize,inner sep=1.5pt},
  fb op/.style={font=\large},
  fb hregion/.style={hard blob,fill=gray!22,draw=gray!60!black,line width=0.6pt,dash pattern=on 2.5pt off 1.5pt}]
  % ======================= left: one K gluon attached to the hard part =======================
  \begin{scope}
    \coordinate (V) at (-0.4,1.55);       % photon vertex
    \coordinate (G) at (1.0,2.3);         % attachment of the K gluon on the struck quark
    \coordinate (E) at (2.05,2.86);       % outgoing struck quark
    \begin{scope}[on background layer]
      \draw[fb hregion,rotate around={28.2:(0.3,1.92)}] (0.3,1.92) ellipse[x radius=1.12,y radius=0.42];
    \end{scope}
    \node[jet blob,ellipse,minimum width=1.9cm,minimum height=0.7cm,inner sep=0pt] (A) at (0,0) {};
    \node[font=\small] at (-0.45,0) {$A$};
    \draw[photon] (-1.75,2.3) -- (V);
    \draw[quark] (-0.4,0.34) -- (V);
    \draw[quark] (V) -- (G);
    \draw[quark] (G) -- (E);
    \draw[longit] (0.6,0.33) to[out=72,in=-108] (G);
    \fill[blue!60!black] ($(G)+(-0.06,-0.19)$) circle[radius=1.7pt];
    \node[vertex] at (V) {}; \node[vertex] at (G) {};
    \path[mom'={$q$},mom sep=0.24cm,mom length=0.45cm] (-1.75,2.3) -- (V);
    \path[mom={$l$},mom sep=0.24cm,mom length=0.45cm] (-0.4,0.34) -- (V);
    \path[mom={$p'-k$},mom sep=0.3cm,mom length=0.42cm,mom pos=0.42] (V) -- (G);
    \path[mom={$p'$},mom sep=0.24cm,mom length=0.42cm] (G) -- (E);
    \path[mom'={$k$},mom sep=0.26cm,mom length=0.45cm,mom pos=0.45] (0.6,0.33) to[out=72,in=-108] (G);
    \node[fb lab,text=gray!35!black] at ($(G)+(0.42,-0.36)$) {$H$};
    \node[fb lab,anchor=north] at (0.08,-0.42) {$k_\mu H^\mu_a=g\,t^a\,H_0$};
    \node[font=\large,anchor=east] at (-1.75,1.2) {$\displaystyle\sum_{H\text{ 内的附着}}$};
  \end{scope}
  \node[fb op] at (2.75,1.4) {$=$};
  % ======================= middle: H_0 =======================
  \begin{scope}[xshift=4.55cm]
    \coordinate (V) at (-0.25,1.65);
    \draw[photon] (-1.3,2.35) -- (V);
    \draw[quark] (-0.25,0.6) -- (V);
    \draw[quark] (V) -- (1.15,2.4);
    \node[vertex] at (V) {};
    \path[mom={$l+k$},mom sep=0.24cm,mom length=0.45cm] (-0.25,0.6) -- (V);
    \path[mom={$p'$},mom sep=0.24cm,mom length=0.42cm] (V) -- (1.15,2.4);
    \node[fb lab] at (0.6,1.42) {$H_0$};
  \end{scope}
  \node[fb op] at (6.15,1.4) {$\times$};
  % ======================= right: the Wilson-line vertex of the PDF operator =======================
  \begin{scope}[xshift=7.75cm]
    \coordinate (O) at (-0.35,1.65);     % operator vertex psi(0)
    \coordinate (W) at (1.75,2.8);       % towards infinity along n_B
    \coordinate (Y) at ($(O)!0.5!(W)$);
    \node[jet blob,ellipse,minimum width=1.9cm,minimum height=0.7cm,inner sep=0pt] (A) at (0,0) {};
    \node[font=\small] at (-0.45,0) {$A$};
    \draw[quark] (-0.35,0.34) -- (O);
    \draw[wilson=0.78] (O) -- (W);
    \node[opvertex] at (O) {};
    \draw[longit] (0.55,0.33) to[out=70,in=-110] (Y);
    \node[vertex] at (Y) {};
    \path[mom'={$k$},mom sep=0.26cm,mom length=0.45cm,mom pos=0.45] (0.55,0.33) to[out=70,in=-110] (Y);
    \node[fb lab,anchor=east] at ($(O)+(-0.12,0.05)$) {$\psi(0)$};
    \node[fb lab,anchor=west] at (W) {$\infty n_B$};
    \node[fb lab,anchor=north west] at ($(Y)+(0.12,-0.06)$) {$\ii g\,t^a n_B^\alpha$};
    \node[fb lab,anchor=south east] at ($($(O)!0.5!(Y)$)+(-0.02,0.12)$) {$\dfrac{-\ii}{n_B\cdot k-\ii0}$};
    \node[fb lab,anchor=north] at (0.4,-0.42) {$W^{(1)}[\infty n_B,0]$：未来指向};
  \end{scope}
  % summary formula
  \node[font=\footnotesize,anchor=north] at (4.4,-1.15) {$H_\mu\,\dfrac{k^\mu n_B^\alpha}{n_B\cdot k-\ii0}=\Big[\ii g\,t^an_B^\alpha\,\dfrac{-\ii}{n_B\cdot k-\ii0}\Big]H_0=W^{(1)}[\infty n_B,0]\,H_0$};
\end{tikzpicture}
\caption{DIS 中 PDF Wilson 线的种子（VF-WL-8，\cref{prop:method-ward-tree}）：一个来自靶共线子图 $A$ 的 K 胶子（蓝色，流入硬部分的动量为 $k$，圆点标出带标量极化 $k_\mu$ 的硬端）对硬部分 $H$ 的所有附着之和（领头阶即附着在末态被击夸克 $p'$ 上），
经 Ward 恒等式 $k_\mu H^\mu_a=g\,t^aH_0$ 化为去掉该胶子、入射动量为 $l+k$ 的硬部分 $H_0$，乘以夸克场 $\psi(0)$ 处出发、沿 $n_B$ 的未来指向 Wilson 线 $W^{(1)}[\infty n_B,0]$ 的单胶子项（双线箭头为路径方向 $0\to\infty n_B$）。
程函传播子中的 $-\ii0$ 继承自被替代的被击夸克传播子：$(p'-k)^2+\ii0\simeq-2p'^-k^++\ii0$ 的极点位于 $k^+$ 上半平面；Drell--Yan 中同样的附着落在入射反夸克上，给出 $+\ii0$。}
\label{fig:disall-seed}
\end{figure}
%%% END fig:disall-seed

\begin{proposition}[从非物理区前向振幅到矩与 $x$ 空间的两条途径 \tagR]\label{prop:disall-two-routes}
（$R_{\mathrm{cond}}$：在所列编号假设之下证明完整；结论的地位不高于其中最弱的假设。）前提：\cref{asm:appN-collins}（在复圆盘 $|\omega|\le\rho<1$ 上一致）、\cref{asm:appN-contour}、\cref{asm:appN-envelope}、\cref{asm:appN-separation}、\cref{asm:appN-cutsum}（用于 (c)）。（本命题只涉及部分子态。）
设部分子靶 $j$ 带有红外正规化（参数 $\lambda$）。在 $\as$ 的固定阶，对前向振幅 $T^{\mu\nu}(\omega)$ 的全部领头幂次图按\cref{def:method-subtraction}构造减除项 $C_RT^{\mu\nu}$，
记余项 $r^{\mu\nu}\equiv T^{\mu\nu}-\sum_RC_RT^{\mu\nu}$，$r_a$（$a=1,2$）为其不变振幅（以\cref{def:disall-breit}的投影提取）。
\begin{enumerate}[label=(\alph*),leftmargin=2em,itemsep=1pt,topsep=2pt]
\item 对 $|\omega|<1$，$T^{\mu\nu}$ 的领头区域只含靶共线模式与硬模式：没有末态，携带 $q$ 的线离壳 $\order{Q^2}$，只附着在靶共线子图上的软线属于共线子图。
因此只用到硬—共线近似子，（取交封闭的）区域集合只涉及一个共线方向；并且对每个实数 $\omega\in[-\rho,\rho]$，$\rho<1$，有 $r^{\mu\nu}(\omega)=\order{\lambda^aL^K}$，
$a>0$ 与 $K$ 不依赖于 $\lambda$，常数对 $\varepsilon\in(0,\varepsilon_0]$ 一致。
\item 若 (a) 的界对复圆盘 $|\omega|\le\rho$ 一致成立，则每个矩都以 $\order{\rho^{-n}\lambda^aL^K}$ 的余项因子化：记 $M^{(2)}_n=\int_0^1\dd x\,x^{n-2}F_2$，
$M^{(1)}_n=\int_0^1\dd x\,x^{n-1}F_1$（$n$ 为偶数），$[\omega^k]f$ 为 $f$ 在 $\omega=0$ 处的第 $k$ 个 Taylor 系数，则
\begin{equation}\label{eq:disall-cauchy}
\begin{gathered}
\Big|M^{(2)}_n-\tfrac14[\omega^{n-1}]\textstyle\sum_RC_RT_2\Big|\le\tfrac14\rho^{1-n}\max_{|\omega|=\rho}|r_2(\omega)|,\\
\Big|M^{(1)}_n-\tfrac14[\omega^{n}]\textstyle\sum_RC_RT_1\Big|\le\tfrac14\rho^{-n}\max_{|\omega|=\rho}|r_1(\omega)|.
\end{gathered}
\end{equation}
仅在实线段上的界不给出 Taylor 系数的界：复圆盘上的一致性是额外的输入。
\item 实线段上的界不控制不连续部分。固定 $x<1$ 处的 $x$ 空间陈述需要在\cref{asm:appN-cutsum}之下对割图求和后的被积函数作减除分析；当 $x\to1$ 时这一分析中的常数不一致（阈值区）。
\end{enumerate}
\uses{def:disall-breit,def:dis-compton,def:method-subtraction,def:method-collinear-approx,def:method-collins-conditions,asm:appN-collins,asm:appN-cutsum,asm:appN-contour,asm:appN-envelope,asm:appN-separation,thm:appN-subtraction,lem:appN-error-table,lem:method-local-accuracy,lem:method-meets,prop:dis-compton-analytic,prop:dis-dispersion}
% errtype: E-N
\end{proposition}

\begin{proof}
(a) 由\cref{eq:disall-unphysical-offshell}，携带 $q$ 的线都离壳 $\ge Q^2(1-\rho)-\order{\lambda Q^2}$，所以每个领头区域是在每条线上指派 $A$（靶共线）或 $H$，携带 $q$ 的线必为 $H$。
区域按逐线偏序比较，交逐线取得（\cref{lem:method-meets}）：钳夹集相交时 $\pi_{R\wedge R'}=\pi_R\cap\pi_{R'}$，交属于交封闭的区域集合；不相交时（例如 $H_R\cap H_{R'}$ 中没有连接两个光子顶点、携带 $q$ 的路径）二者没有公共下界，(CS2) 直接要求 $C_R\Gamma\in\mathfrak S(U_{R'})$。
$T_R$ 在 $H_R$ 中作 $k\to\pi_+(k)=(k^+,0,\mathbf 0_T)$ 并在界面插入 Dirac 投影 $\frac12\gamma^-\gamma^+$ 或胶子投影（\cref{def:method-collinear-approx}，只用误差表中共线近似的各行）；它不引入新分母，
被投影的硬分母在 $U_R$ 上由\cref{asm:appN-separation}(i) 满足 $|l^2-M_l^2|\ge c\,t_R^2$，非物理区又没有末态割线，因此 $C_RT^{\mu\nu}$ 在\cref{asm:appN-contour}的围道上没有新奇点，不出现软近似分母的 $\ii0$ 相容性问题。
极小区域的 (CS1) 由\cref{lem:method-local-accuracy}给出（前提\cref{asm:appN-envelope,asm:appN-separation}，误差各行见\cref{lem:appN-error-table}），其余的 (CS1)(CS2)(CS3) 是\cref{asm:appN-collins}的内容，
于是\cref{thm:appN-subtraction}（前提\cref{asm:appN-contour}(a)）给出 $\int|r|\le C\lambda^aL^K$，对 $\varepsilon\in(0,\varepsilon_0]$ 一致；\cref{asm:appN-collins}在圆盘上一致，故常数对 $\omega\in[-\rho,\rho]$ 一致。

(b) 在固定阶，$T_{1,2}(\omega)$ 在去掉实轴上 $|\omega|\ge1$ 两条割线的平面上解析（\cref{prop:dis-compton-analytic}）；$C_RT^{\mu\nu}$ 只通过入射动量为 $\xi p$（$|\xi|\le1$）的被投影硬因子依赖于 $\omega$，
其阈值在 $(\xi p\pm q)^2=0$ 即 $\xi\omega=\pm1$，故 $r_a$ 在 $|\omega|<1$ 中解析，Cauchy 估计给出 $|[\omega^k]r_a|\le\rho^{-k}\max_{|\omega|=\rho}|r_a|$。由色散关系（\cref{prop:dis-dispersion}）
$T_2=\sum_{n\ \text{偶}}4\omega^{n-1}M^{(2)}_n$，$T_1-T_1(0)=\sum_{n\ \text{偶}}4\omega^nM^{(1)}_n$，即得\cref{eq:disall-cauchy}。反之，$\lambda^a\sin(\omega/\lambda)$ 在实轴上不超过 $\lambda^a$，
其奇数阶 Taylor 系数却为 $\pm\lambda^{a-k}/k!$，对 $k>a$ 随 $\lambda\to0$ 发散。

(c) 不连续部分的逐点大小不受 $|\omega|<1$ 上取值的控制：若某余项的不连续部分为 $\sin(x/\lambda^a)$（$0\le x\le1$），分部积分给出 $|\int_0^1\dd x\,x^n\sin(x/\lambda^a)|\le2\lambda^a$（一切 $n\ge0$），
由色散关系它在 $|\omega|\le\rho$ 上对应的前向振幅不超过 $8\rho\lambda^a/(1-\rho^2)$，而它本身逐点为 $\order{1}$（同类反例见\cref{rem:dis-nonuniform}）。所以固定 $x$ 须直接处理割图：
在\cref{asm:appN-cutsum}之下，每个前向图 $G$ 的割图之和在公共区域 $\Omega_G$ 上有公共被积函数 $F_G$（光锥序表示），近似子与割图求和对易，(GA1)–(GA4)、(CS1)–(CS3) 对 $F_G$ 成立，
\cref{thm:appN-subtraction}便可在固定 $x$ 处应用于 $F_G$；其中不变质量为 $Q^2(\xi-x)/x$ 的末态整体并入硬子图，内部的末态共线与软钳夹在割图求和后相消。
当 $\xi\to x$ 时这一不变质量趋于零，流喷注成为真正的共线子图；对 $1-x\lesssim\lambda^b$（某个 $b>0$）不存在使流喷注为硬的区域，常数随 $x\to1$ 增长：这一端点正是加号分布与阈值区的来源（\cref{ch:thr}）。
\end{proof}

\noindent
(b) 是“矩的途径”：不需要 OPE，但需要复 $\omega$ 上的一致界，只给出分布意义下的结论；(c) 是“$x$ 空间的途径”，代价是割图求和的公共表示\cref{asm:appN-cutsum}，主定理\cref{thm:disall-factorization}走这一条路。
